\section{Problemstilling}
\label{sec:problemstilling}

HDO har to videomotorer i drift, hver med egne API-er og
integrasjonsmetoder. Eksisterende løsninger er tett koblet til én backend
om gangen, noe som vanskeliggjør leverandørbytte, utvidelse med nye
tjenester og standardisering av funksjonalitet på tvers av motorene. Siden
systemet brukes i akuttmedisinsk sammenheng, må endringer samtidig ivareta
krav til sikkerhet, tilgjengelighet, stabilitet og responstid.

% Dette prosjektet utviklet og evaluerte en plugin-basert mellomvare, VideoAPI, som standardiserer tilgang til ulike videomotorer. Løsningen ble implementert som proof-of-concept med støtte for AntMedia Server og NHN Video.

Det sentrale spørsmålet oppgaven besvarer er:

\begin{quote}
\textit{I hvilken grad kan en plugin-basert mellomvarearkitektur
standardisere tilgang til AntMedia Server og NHN Video gjennom ett felles
API, samtidig som backend-spesifikk logikk isoleres fra kjernelogikken i
VideoAPI?}
\end{quote}

Problemstillingen vurderes langs fire dimensjoner. For det første vurderes
om et felles API dekker kjerneoperasjonene i kravspesifikasjonen. For det
andre vurderes om forskjeller mellom AntMedia Server og NHN Video kan
isoleres i provider-laget. For det tredje vurderes om sikkerhets-,
infrastruktur- og observability-kravene oppfylles. For det fjerde diskuteres
begrensningene som oppstår når to videomotorer har ulike domenemodeller.

Oppgavens arbeidshypotese er:

\begin{quote}
\textit{En modulær mellomvare med provider-basert arkitektur kan redusere
klientenes integrasjonskompleksitet ved å samle tilgang til heterogene
videomotorer bak ett felles API, men backend-spesifikke forskjeller må
fortsatt håndteres eksplisitt i provider-laget.}
\end{quote}